% Data Protection Notice, officer variant.
%
% A SEPARATE FILE, AND NOT ONLY FOR WORDING. The intern notice lists what is collected for
% an INTERN, and the officer class collects materially less: person-classes.yaml forbids
% `education` and `statutory` and does not permit `emergency_contact`. A notice that
% overstates what is held is a defect in the notice -- the whole purpose of the document is
% to be an accurate account of the processing.
%
% So section 1 below is derived from what the officer class actually permits, and section 2
% from what those fields are actually for. Change person-classes.yaml and this needs
% changing with it; nothing checks that, and DECISION-DEBT records it.
% adapted-from: templates/policy/data-protection-notice.tex@b92d8e1ca4e13146d61f494331dca8290a4b20e979fd8e59dc83ab04f84f101a
%
% NOT A SHARED FILE, AND NOT A COPY THAT DRIFTS SILENTLY. Of the sections in this pair,
% more than half differ in substance -- two of them inverted -- so there is no shared
% clause to extract. What there is instead is a risk: the intern document is amended and
% this one is not, and nothing notices.
%
% The digest above is of the file this was adapted from. `cargo test --test repository`
% recomputes it and fails when it moves, naming this file. That is the whole mechanism:
% it does not compare content -- the two are SUPPOSED to differ -- it makes it impossible
% to change one half without being told to review the other.
%
% When the intern document legitimately changes: review this one, make the matching
% change or decide it does not apply, then update the digest in the same commit.

\documentclass{paridyum}
\begin{document}

\paridyumletterhead
\paridyumdateref{\dkDocDate}{\dkDocRef}
\paridyumtitle{Data Protection Notice (Officer)}

This notice explains what personal data \textbf{\dkOrgName} collects about
\textbf{\dkPIdentityLegalName} in connection with the office they hold, why it is collected, how
long it is kept, and what rights the Officer has.

\section*{1. What Paridyum Collects}
\begin{itemize}[leftmargin=*, noitemsep]
    \item identity and contact details --- name, preferred name, date of birth and
          email address;
    \item details of the office --- the office held and the date the appointment took
          effect;
    \item professional details voluntarily provided --- such as a portfolio,
          repository or professional profile;
    \item device and access details --- the device used for Paridyum work, accounts
          issued, and access granted and removed;
    \item equipment issued by Paridyum, where any is; and
    \item records of documents issued to and signed by the Officer.
\end{itemize}

Paridyum does not require and does not seek health records, financial account details,
government identity numbers, biometric data, or educational records for the holding of
office. No remuneration is paid for the office and no payroll or tax data is collected.

Date of birth is collected for one reason: the Digital Personal Data Protection Act
treats a person under eighteen differently, and Paridyum cannot apply that rule without
being able to tell.

\section*{2. Why Paridyum Processes It}
Personal data is processed for purposes connected with the office, namely:
\begin{itemize}[leftmargin=*, noitemsep]
    \item recording who holds which office at Paridyum, and from when;
    \item issuing and managing system access, and keeping systems and information
          secure;
    \item preparing the documents that form the officer record, including the
          Non-Disclosure Agreement and the Intellectual Property Assignment;
    \item maintaining the register of who signed what, which must stay verifiable
          after the office is vacated; and
    \item meeting Paridyum's legal and record-keeping obligations.
\end{itemize}

Personal data is not used for marketing, is not sold, and is not used to train
external systems.

\section*{3. Who It Is Shared With}
Personal data is available to those at Paridyum who need it to administer the office.
It is shared outside Paridyum only where the Officer asks for it, or where disclosure is
required by law.

Two facts about an officer are PUBLISHED rather than merely held: their name and the
office they hold appear on Paridyum's public website, together with a summary the Officer
may review and ask to change. Nothing else in this notice is published.

\section*{4. How Long It Is Kept}
The officer record is kept for as long as the office is held and for the retention period
that follows, as set out in Paridyum's retention schedule. Two things are kept longer,
and deliberately: the signed agreements, and the register entry recording that they were
issued. Their purpose is to remain provable --- an intellectual-property assignment has no
end, and a register that cannot say who was an officer in a given year cannot verify a
document issued in it. Personal data no longer needed for any of these purposes is
deleted.

\section*{5. The Officer's Rights}
The Officer may:
\begin{itemize}[leftmargin=*, noitemsep]
    \item ask what personal data Paridyum holds about them and how it is processed;
    \item ask for inaccurate or incomplete data to be corrected or completed;
    \item ask for personal data to be erased where it is no longer required for the
          purposes above and Paridyum is not required to keep it; and
    \item raise a concern about how their personal data has been handled.
\end{itemize}

\section*{6. How to Exercise Them}
Write to \textbf{\dkOrgEmail}. Paridyum will respond within a reasonable period. If the
Officer is not satisfied with the response, they may complain to the Data Protection
Board of India. Holding office at Paridyum does not reduce these rights, and an officer
exercising them is treated no differently from anybody else.

\section*{7. Keeping Data Accurate}
The Officer should tell Paridyum if any of the information above changes while they hold
office, so that the record stays accurate.

\section*{8. Acknowledgement}
By signing below, the Officer confirms that they have received and read this notice.

\vspace{1.5em}
\paridyumsignblock

\vspace{1.5em}
\paridyumcountersign{Officer}{\dkPIdentityLegalName}

\end{document}
